\documentclass[11pt,a4paper]{article}

\usepackage[margin=1.1in]{geometry}
\usepackage[T1]{fontenc}
\usepackage[utf8]{inputenc}
\usepackage{lmodern}
\usepackage{amsmath,amssymb,amsthm,mathtools,mathrsfs}
\usepackage{enumitem}
\usepackage{microtype}
\usepackage[colorlinks=true,linkcolor=blue,citecolor=blue,urlcolor=blue]{hyperref}

\newtheorem{theorem}{Theorem}[section]
\newtheorem{proposition}[theorem]{Proposition}
\newtheorem{lemma}[theorem]{Lemma}
\newtheorem{corollary}[theorem]{Corollary}
\theoremstyle{definition}
\newtheorem{definition}[theorem]{Definition}
\theoremstyle{remark}
\newtheorem{remark}[theorem]{Remark}

\DeclareMathOperator{\Fix}{Fix}
\DeclareMathOperator{\Ann}{Ann}
\DeclareMathOperator{\Spec}{Spec}
\DeclareMathOperator{\ord}{ord}
\DeclareMathOperator{\Idl}{Idl}
\DeclareMathOperator{\Tail}{Tail}
\DeclareMathOperator{\supp}{supp}
\DeclareMathOperator{\width}{width}
\DeclareMathOperator{\im}{im}
\DeclareMathOperator{\coker}{coker}

\newcommand{\ZZ}{\mathbb Z}
\newcommand{\CC}{\mathbb C}
\newcommand{\RR}{\mathbb R}
\newcommand{\PP}{\mathbb P}
\newcommand{\BB}{\mathbf B}
\newcommand{\NN}{\mathbb N}
\newcommand{\cO}{\mathcal O}
\newcommand{\cP}{\mathcal P}
\newcommand{\cV}{\mathcal V}
\newcommand{\cA}{\mathcal A}
\newcommand{\cQ}{\mathcal Q}
\newcommand{\cK}{\mathcal K}
\newcommand{\cI}{\mathcal I}
\newcommand{\cB}{\mathcal B}
\newcommand{\cC}{\mathcal C}
\newcommand{\taup}{\tau}
\newcommand{\into}{\hookrightarrow}
\newcommand{\onto}{\twoheadrightarrow}

\title{Boolean Cube Shadows, Jet-Ideal Chains,\
Compactified Tail Sheaves, and Packet Descent\
for Zeta Zero Divisors}
\author{}
\date{May 2026}

\begin{document}
\maketitle

\begin{abstract}
Let
\[
S=G(\ZZ)=\ZZ\sqcup\{\taup\},\qquad A:=\{0,\taup\}\cong\BB,
\qquad F_1(s):=\zeta(s)-1.
\]
At every zero $\rho$ of the Riemann zeta function one has $F_1(\rho)=-1$, and multiplication by $-1$ fixes exactly the Boolean pair $A$.  The pointwise split-zero shadow is therefore constant.  This paper constructs the first genuinely global and genuinely nonreduced refinements of that observation and proves the precise relations among them.

For an effective discrete divisor $D=\sum_x m_x[x]$ on a Riemann surface, we define the Boolean multiplicity shadow
\[
\cV_D(U)=\prod_{x\in U\cap |D|}A^{m_x}
\]
and the reduced shadow
\[
\cA_D(U)=\prod_{x\in U\cap |D|}A.
\]
We prove that $\cA_D$ is a canonical split retract of $\cV_D$, and that $\cV_D\cong\cA_D$ holds in sheaves of $A$-semimodules if and only if every multiplicity equals $1$.

We then introduce the coordinate-free classical jet thickening
\[
Q_{x,m_x}:=\cO_{X,x}/\mathfrak m_x^{m_x}
\]
and prove that its ideal lattice is canonically the chain $C_{m_x}$ of terminal segments in $[m_x]$.  The Boolean cube $B_{m_x}=\cP([m_x])$ contains $C_{m_x}$ as a canonical sublattice, and we prove the explicit biadjoint relation
\[
L_{m_x}\dashv i_{m_x}\dashv R_{m_x}
\]
inside finite posets, where $i_{m_x}:C_{m_x}\into B_{m_x}$ is the terminal-segment inclusion.  Thus the classical nonreduced divisor and the Boolean cube shadow are not identified; instead they are related by a reflective-coreflective comparison theorem.

For a holomorphic function $f$ with divisor $D$, the class of $f$ in the $(m_x+1)$-jet algebra
\[
K_{x,m_x}:=\cO_{X,x}/\mathfrak m_x^{m_x+1}
\]
is proved to generate the minimal nonzero ideal, giving the requested derivative-enhanced stalk.  On $\CC$, this recovers the exact coefficient $f^{(m_x)}(x)/m_x!$.

Finally, after compactification to $\PP^1(\CC)$, we compute the direct image and extension-by-zero of these shadows and identify the new fiber at $\infty$ as the tail object
\[
\Tail\bigl((M_x)_{x\in |D|}\bigr)=\prod_x M_x\Big/\sim
\]
(modulo finite modification).  Applied to zeta, this produces an honest asymptotic Boolean tail sheaf.  For the completed zeta function $\xi$, we also prove descent of the nontrivial-zero shadow to the Klein-four orbit space and obtain a canonical packet sheaf.  On the trivial-zero sector we prove the exact weight formula
\[
|\zeta'(-2n)|=\frac{(2n)!}{2(2\pi)^{2n}}\,\zeta(2n+1),
\]
which calibrates the unique nonempty face of the trivial Boolean stalk.
\end{abstract}

\tableofcontents

\section{Introduction}

The pointwise identity
\[
\zeta(\rho)=0\implies F_1(\rho)=\zeta(\rho)-1=-1
\]
by itself does almost nothing: every zeta zero lands on the same supported unit of the globalization semiring $G(\ZZ)$.  The mathematical problem is therefore not pointwise identification but how to organize the whole zero divisor into a meaningful geometric object.

There are at least three different objects one can attach to a zero of multiplicity $m$.

\begin{enumerate}[label=(\roman*)]
\item The reduced Boolean stalk $A$, which remembers only whether the zero is present.
\item The Boolean cube $A^m\cong\cP([m])$, which remembers $m$ abstract layers.
\item The classical nonreduced local algebra $Q_m\cong \CC[t]/(t^m)$, whose ideal lattice is the chain of powers of $t$.
\end{enumerate}

The second and third objects should not be conflated.  They are not isomorphic in any reasonable ordered category once $m\ge2$.  But they are not unrelated either.  The key theorem of this paper is that the ideal chain of the classical jet thickening sits inside the Boolean cube as the terminal-segment chain, and this inclusion admits both a left adjoint and a right adjoint.  This gives a precise positive replacement for the vague statement that the Boolean cube is ``like'' a nonreduced thickening.

The second theme is globality.  Because zero sets of holomorphic functions on one-dimensional complex manifolds are discrete, any divisor-shadow sheaf on the open surface has vanishing higher cohomology.  That observation is correct, but it is not the end of the story.  After compactification to $\PP^1(\CC)$, the direct image and extension-by-zero differ by an explicit tail object at $\infty$, namely the quotient of the full product by finite modification.  This tail fiber is the first genuinely global invariant carried by the split-zero shadow in the present setting.

The third theme is symmetry.  The completed zeta function
\[
\xi(s)=\frac12 s(s-1)\pi^{-s/2}\Gamma\!\left(\frac s2\right)\zeta(s)
\]
is invariant under $s\mapsto 1-s$ and under complex conjugation.  Those symmetries organize the nontrivial zeros into Klein-four packets.  We prove that the multiplicity shadow and the jet-ideal chain both descend to the orbit space as actual packet sheaves, not merely as rhetoric.

The background frameworks relevant here are standard: sheaves and stalks on complex manifolds \cite{Lee}, affine and local constructions in algebraic geometry \cite{GrothendieckEGAII,Vakil}, idempotent and semiring geometry \cite{GaubertKatz,Gunawardena,DrosteKuichVogler}, and the embedding of semirings into blueprints \cite{Lorscheid}.  The paper uses these only as grounding.  Every statement needed in the sequel is proved directly.

\section{Split-root fixed loci in globalization semirings}

\begin{definition}
Let $R$ be a commutative ring with identity.  Its \emph{globalization semiring} is
\[
G(R):=R\sqcup\{\taup\},
\]
with operations
\[
a\oplus b:=a+b,\qquad a\oplus\taup=\taup\oplus a:=a,\qquad \taup\oplus\taup:=\taup,
\]
and
\[
a\odot b:=ab,\qquad a\odot\taup=\taup\odot a:=\taup,\qquad \taup\odot\taup:=\taup.
\]
Thus $\taup$ is the additive identity and multiplicative absorber, while $1\in R\subset G(R)$ is the multiplicative identity.
\end{definition}

\begin{definition}
For $u\in R^\times\subset G(R)^\times$, define its \emph{split-root shadow} to be
\[
\Fix_{G(R)}(u):=\{x\in G(R):u\odot x=x\}.
\]
\end{definition}

\begin{lemma}\label{lem:units-fix}
For every commutative ring $R$,
\[
G(R)^\times=R^\times.
\]
Moreover, for every $u\in R^\times$,
\[
\Fix_{G(R)}(u)=\{\taup\}\cup\Ann_R(u-1).
\]
\end{lemma}

\begin{proof}
If $x\in G(R)$ is invertible, say $x\odot y=1$, then $x\neq \taup$ because $\taup\odot y=\taup\neq1$.  Likewise $y\neq\taup$.  Hence $x,y\in R$ and $xy=1$ inside $R$, so $x\in R^\times$.  The converse is immediate.

Now let $u\in R^\times$.  Since $u\odot\taup=\taup$, the absorber is always fixed.  For $r\in R$ one has
\[
u\odot r=r \iff ur=r \iff (u-1)r=0,
\]
which proves the formula.
\end{proof}

\begin{proposition}\label{prop:split-root-domain}
Let $R$ be an integral domain and let $u\in R^\times$ satisfy $u\neq1$.  Then
\[
\Fix_{G(R)}(u)=\{0,\taup\}.
\]
In particular, if $\mathcal O_K$ is the ring of integers of a number field and $\xi\in\mathcal O_K^\times$ is a nontrivial root of unity, then
\[
\Fix_{G(\mathcal O_K)}(\xi)=\{0,\taup\}.
\]
\end{proposition}

\begin{proof}
By Lemma \ref{lem:units-fix},
\[
\Fix_{G(R)}(u)=\{\taup\}\cup\Ann_R(u-1).
\]
Because $R$ is a domain and $u-1\neq0$, multiplication by $u-1$ is injective.  Hence $\Ann_R(u-1)=\{0\}$.
\end{proof}

\begin{proposition}\label{prop:boolean-pair}
For every ring $R$, the subset
\[
A_R:=\{0,\taup\}\subset G(R)
\]
is a subsemiring canonically isomorphic to the Boolean semiring $\BB=\{0,1\}$.  Explicitly,
\[
\taup\longleftrightarrow 0\in\BB,
\qquad
0\longleftrightarrow 1\in\BB.
\]
\end{proposition}

\begin{proof}
Inside $A_R$ the operation tables are
\[
0\oplus0=0,\qquad 0\oplus\taup=0,\qquad \taup\oplus\taup=\taup,
\]
and
\[
0\odot0=0,\qquad 0\odot\taup=\taup,\qquad \taup\odot\taup=\taup.
\]
These are exactly the Boolean rules after the displayed identification.
\end{proof}

\begin{corollary}\label{cor:zeta-pointwise}
Let $S=G(\ZZ)$ and $F_1(s)=\zeta(s)-1$.  If $\rho$ is any zero of $\zeta$, then
\[
F_1(\rho)=-1,
\qquad
\Fix_S(F_1(\rho))=\{0,\taup\}\cong\BB.
\]
\end{corollary}

\begin{proof}
The first identity is immediate from $\zeta(\rho)=0$.  The second is Proposition \ref{prop:split-root-domain} with $R=\ZZ$ and $u=-1$.
\end{proof}

\begin{remark}
The idempotent role of the adjoined element $\taup$ is formally parallel to the bottom element in max-plus and tropical settings \cite{GaubertKatz,Gunawardena}.  The zeta fixed-locus phenomenon is therefore naturally an idempotent-support phenomenon rather than an ordinary ring-theoretic one.
\end{remark}

\section{Boolean divisor shadows on a Riemann surface}

Throughout this section, $X$ is a connected Riemann surface and
\[
D=\sum_{x\in Z} m_x[x]
\]
is an effective divisor with discrete support $Z\subset X$.  We write $A:=\{0,\taup\}\cong\BB$.

\begin{definition}
The \emph{Boolean multiplicity shadow} of $D$ is the sheaf of $A$-semimodules
\[
\cV_D(U):=\prod_{x\in U\cap Z}A^{m_x}
\]
on open sets $U\subset X$, with restriction maps given by coordinate projection.  The \emph{reduced Boolean shadow} is
\[
\cA_D(U):=\prod_{x\in U\cap Z}A.
\]
\end{definition}

\begin{lemma}\label{lem:product-sheaf}
Let $Z\subset X$ be closed and discrete, and let $(M_x)_{x\in Z}$ be a family of sets, posets, semimodules, or vector spaces.  Then
\[
\mathcal F_M(U):=\prod_{x\in U\cap Z}M_x
\]
defines a sheaf on $X$.
\end{lemma}

\begin{proof}
Let $U=\bigcup_{i\in I}U_i$ be an open cover.  If two sections of $\mathcal F_M(U)$ agree on each $U_i$, then their $x$-coordinates agree for every $x\in U\cap Z$ because some $U_i$ contains $x$.  Hence the sections are equal.

Conversely, suppose $s_i\in\mathcal F_M(U_i)$ are compatible.  For each $x\in U\cap Z$, choose $i$ with $x\in U_i$ and define the $x$-coordinate of the glued section to be the $x$-coordinate of $s_i$.  Compatibility on overlaps makes this independent of the choice of $i$.  This produces the unique glued section on $U$.
\end{proof}

\begin{proposition}\label{prop:shadow-stalks}
The sheaves $\cV_D$ and $\cA_D$ are well defined.  Their stalks at $x\in Z$ are
\[
(\cV_D)_x\cong A^{m_x},
\qquad
(\cA_D)_x\cong A,
\]
and their stalks away from $Z$ are one-point stalks.
\end{proposition}

\begin{proof}
This is Lemma \ref{lem:product-sheaf}.  Since $Z$ is discrete, for each $x\in Z$ there is an open neighborhood $U$ with $U\cap Z=\{x\}$.  For every smaller neighborhood $V\subset U$ of $x$ one has
\[
\cV_D(V)=A^{m_x},\qquad \cA_D(V)=A,
\]
with identity restrictions, so the direct limits defining the stalks are $A^{m_x}$ and $A$.  If $x\notin Z$, choose $U$ with $U\cap Z=\varnothing$; then all sections on smaller neighborhoods are singleton products.
\end{proof}

\begin{proposition}\label{prop:boolean-cube}
For every integer $m\ge1$, the free rank-$m$ $A$-semimodule $A^m$ is canonically isomorphic to the Boolean cube
\[
B_m:=\cP([m]),\qquad [m]:=\{1,\dots,m\},
\]
by the rule
\[
(a_1,\dots,a_m)\longmapsto \{j:a_j=0\}.
\]
Under this identification the semimodule addition is union of subsets.
\end{proposition}

\begin{proof}
Because $A\cong\BB$, coordinatewise addition on $A^m$ is Boolean OR.  Therefore
\[
\{j:(a_j\oplus b_j)=0\}=\{j:a_j=0\text{ or }b_j=0\},
\]
which is exactly the union of the two corresponding subsets.  Bijectivity is obvious.
\end{proof}

\begin{proposition}\label{prop:retract-shadow}
For each $m\ge1$, define
\[
\Sigma_m:A^m\to A,
\qquad
\Sigma_m(a_1,\dots,a_m):=a_1\oplus\cdots\oplus a_m,
\]
and
\[
\Delta_m:A\to A^m,
\qquad
\Delta_m(a):=(a,\dots,a).
\]
Then $\Sigma_m$ and $\Delta_m$ are $A$-semimodule maps and
\[
\Sigma_m\circ\Delta_m=\operatorname{id}_A.
\]
Consequently there are canonical sheaf morphisms
\[
\Sigma_D:\cV_D\to\cA_D,
\qquad
\Delta_D:\cA_D\to\cV_D,
\qquad
\Sigma_D\circ\Delta_D=\operatorname{id}_{\cA_D}.
\]
\end{proposition}

\begin{proof}
If $a=\taup$, then $\Sigma_m(\Delta_m(a))=\taup$.  If $a=0$, then idempotence in $A$ gives $\Sigma_m(\Delta_m(a))=0$.  Everything globalizes coordinatewise over the divisor support.
\end{proof}

\begin{theorem}\label{thm:reduced-vs-multiplicity}
The following are equivalent.
\begin{enumerate}[label=(\roman*)]
\item $\cV_D\cong\cA_D$ in sheaves of $A$-semimodules.
\item $m_x=1$ for every $x\in Z$.
\end{enumerate}
When these conditions fail, $\cA_D$ is nevertheless a canonical split retract of $\cV_D$.
\end{theorem}

\begin{proof}
If every $m_x=1$, then the definitions give $\cV_D=\cA_D$.

Conversely, suppose $\cV_D\cong\cA_D$.  Then the stalks are isomorphic at every $x\in Z$, so
\[
A^{m_x}\cong A
\]
in $A$-semimodules.  Any such isomorphism is in particular a bijection of underlying sets, so
\[
2^{m_x}=|A^{m_x}|=|A|=2.
\]
Thus $m_x=1$ for all $x$.  The split-retract statement is Proposition \ref{prop:retract-shadow}.
\end{proof}

\begin{remark}
The sheaf $(Z,\underline{\BB})$ on the discrete support $Z$ is a $0$-dimensional semiring scheme because $\Spec(\BB)$ has one point.  By the embedding of semirings into blueprints, the Boolean shadow also lives naturally inside blueprint geometry \cite{Lorscheid}.  We do not use that machinery below, but the shadow is an honest semiring object rather than an analogy.
\end{remark}

\section{Jet thickenings and the chain-cube comparison theorem}

\subsection{Classical local Artin rings}

\begin{definition}
For $x\in X$ and $m\ge1$, let $\mathfrak m_x\subset\cO_{X,x}$ be the maximal ideal.  Define the local Artin rings
\[
Q_{x,m}:=\cO_{X,x}/\mathfrak m_x^m,
\qquad
K_{x,m}:=\cO_{X,x}/\mathfrak m_x^{m+1}.
\]
For the divisor $D=\sum_{x\in Z}m_x[x]$, set
\[
Q_x:=Q_{x,m_x},\qquad K_x:=K_{x,m_x}.
\]
\end{definition}

\begin{lemma}\label{lem:dvr-local}
Let $x\in X$.  Then there exists a holomorphic coordinate $t$ centered at $x$ such that
\[
\cO_{X,x}\cong \CC\{t\},\qquad \mathfrak m_x=(t).
\]
Consequently
\[
Q_{x,m}\cong \CC\{t\}/(t^m)\cong \CC[t]/(t^m),
\qquad
K_{x,m}\cong \CC\{t\}/(t^{m+1})\cong \CC[t]/(t^{m+1}).
\]
\end{lemma}

\begin{proof}
The existence of a local holomorphic coordinate is the defining local property of a Riemann surface \cite[Chapter~1]{Lee}.  In such a coordinate, the germ ring is the convergent power-series ring $\CC\{t\}$ and the maximal ideal is generated by $t$.  Modulo $(t^m)$ every convergent power series is represented by its degree $<m$ truncation, which yields the polynomial description.
\end{proof}

\begin{proposition}\label{prop:ideal-chain-local}
For every $m\ge1$, the ideals of $Q_{x,m}$ are exactly
\[
I_k:=(t^k)/(t^m),\qquad k=0,1,\dots,m,
\]
where $t$ is any centered local coordinate.  In particular,
\[
Q_{x,m}=I_0\supset I_1\supset\cdots\supset I_m=0
\]
is a chain, independent of the chosen coordinate.
\end{proposition}

\begin{proof}
By Lemma \ref{lem:dvr-local}, $Q_{x,m}\cong\CC\{t\}/(t^m)$.  Ideals of the quotient correspond to ideals of $\CC\{t\}$ containing $(t^m)$.  Since $\CC\{t\}$ is a discrete valuation ring, every ideal is $(t^k)$ for a unique $k\ge0$.  Those containing $(t^m)$ are precisely $(t^k)$ with $0\le k\le m$, giving the displayed chain.  Because the ideal lattice of a ring is intrinsic, the description is coordinate-independent.
\end{proof}

\begin{definition}
For $m\ge1$, define the terminal-segment chain
\[
C_m:=\{T_r:1\le r\le m+1\}\subset \cP([m]),
\]
where
\[
T_r:=\{r,r+1,\dots,m\}\quad (1\le r\le m),
\qquad
T_{m+1}:=\varnothing.
\]
We write $B_m:=\cP([m])$ for the full Boolean cube.
\end{definition}

\begin{corollary}\label{cor:ideal-chain-identification}
For every $x\in X$ and $m\ge1$, the ideal lattice $\Idl(Q_{x,m})$ is canonically isomorphic to the chain $C_m$ via
\[
I_k=(t^k)/(t^m)\longmapsto T_{k+1}.
\]
\end{corollary}

\begin{proof}
Proposition \ref{prop:ideal-chain-local} gives all ideals and their inclusion relations.  The displayed bijection preserves order because
\[
I_k\subseteq I_\ell\iff k\ge \ell
\iff T_{k+1}\subseteq T_{\ell+1}.
\]
\end{proof}

\subsection{The chain-cube comparison theorem}

\begin{definition}
Define the inclusion
\[
i_m:C_m\into B_m
\]
by viewing each terminal segment as a subset of $[m]$.  Define monotone maps
\[
L_m:B_m\to C_m,
\qquad
L_m(I):=\begin{cases}
\varnothing,& I=\varnothing,\\
\{\min I,\min I+1,\dots,m\},& I\neq\varnothing,
\end{cases}
\]
and
\[
R_m:B_m\to C_m,
\qquad
R_m(I):=\{j\in [m]:\{j,j+1,\dots,m\}\subseteq I\}.
\]
\end{definition}

\begin{lemma}\label{lem:LRtails}
For every $I\in B_m$, the sets $L_m(I)$ and $R_m(I)$ lie in $C_m$.  Moreover, $L_m(I)$ is the least terminal segment containing $I$, and $R_m(I)$ is the greatest terminal segment contained in $I$.
\end{lemma}

\begin{proof}
If $I=\varnothing$, then $L_m(I)=\varnothing=T_{m+1}\in C_m$.  If $I\neq\varnothing$, then $L_m(I)=T_{\min I}\in C_m$ and by construction it is the least terminal segment containing $I$.

For $R_m(I)$, set
\[
r:=\min\{j\in[m]:T_j\subseteq I\}
\]
if such a $j$ exists, and $r:=m+1$ otherwise.  Then the defining formula gives $R_m(I)=T_r\in C_m$.  By minimality of $r$, this is the greatest terminal segment contained in $I$.
\end{proof}

\begin{theorem}[Chain-cube comparison]\label{thm:chain-cube}
For every $m\ge1$ one has
\[
L_m\dashv i_m\dashv R_m
\]
in the category of finite posets.  Explicitly, for $I\in B_m$ and $T\in C_m$,
\[
L_m(I)\subseteq T \iff I\subseteq i_m(T),
\]
and
\[
i_m(T)\subseteq I \iff T\subseteq R_m(I).
\]
Hence $C_m$ is both reflective and coreflective in $B_m$.
\end{theorem}

\begin{proof}
By Lemma \ref{lem:LRtails}, $L_m(I)$ is the least terminal segment containing $I$.  Therefore
\[
L_m(I)\subseteq T\iff I\subseteq T=i_m(T).
\]
This proves $L_m\dashv i_m$.

Similarly, $R_m(I)$ is the greatest terminal segment contained in $I$.  Therefore
\[
i_m(T)=T\subseteq I\iff T\subseteq R_m(I),
\]
which proves $i_m\dashv R_m$.
\end{proof}

\begin{corollary}\label{cor:not-iso-chain-cube}
If $m\ge2$, then $C_m$ and $B_m$ are not isomorphic in finite posets, hence not in finite distributive lattices.  Nevertheless they are related by the biadjoint comparison of Theorem \ref{thm:chain-cube}.
\end{corollary}

\begin{proof}
The chain $C_m$ has width $1$.  The Boolean cube $B_m$ has width at least $2$ because $\{1\}$ and $\{2\}$ are incomparable.  Poset isomorphisms preserve width, so no isomorphism exists when $m\ge2$.  The positive relation is Theorem \ref{thm:chain-cube}.
\end{proof}

\subsection{Sheafified chain-cube comparison}

\begin{definition}
Define sheaves of posets on $X$ by
\[
\cI_D(U):=\prod_{x\in U\cap Z}\Idl(Q_x),
\qquad
\cC_D(U):=\prod_{x\in U\cap Z}C_{m_x},
\qquad
\cB_D(U):=\prod_{x\in U\cap Z}B_{m_x},
\]
with coordinatewise order and restriction by projection.
\end{definition}

\begin{proposition}\label{prop:sheaves-posets}
The three assignments $\cI_D$, $\cC_D$, and $\cB_D$ are sheaves of finite posets.  There is a canonical isomorphism
\[
\cI_D\cong \cC_D.
\]
Moreover, under the identification of Proposition \ref{prop:boolean-cube}, the underlying ordered sheaf of the $A$-semimodule sheaf $\cV_D$ is canonically isomorphic to $\cB_D$.
\end{proposition}

\begin{proof}
This is Lemma \ref{lem:product-sheaf} applied to families of finite posets, together with Corollary \ref{cor:ideal-chain-identification} and Proposition \ref{prop:boolean-cube} on each stalk.
\end{proof}

\begin{theorem}\label{thm:sheafified-chain-cube}
There are canonical sheaf morphisms of posets
\[
L_D:\cB_D\to\cC_D,
\qquad
i_D:\cC_D\to\cB_D,
\qquad
R_D:\cB_D\to\cC_D,
\]
obtained coordinatewise from $L_{m_x}$, $i_{m_x}$, and $R_{m_x}$.  Stalkwise one has
\[
(L_D)_x\dashv (i_D)_x\dashv (R_D)_x.
\]
Thus the classical jet-ideal chain sheaf is simultaneously a reflective and a coreflective subobject of the Boolean multiplicity-cube sheaf in sheaves of finite posets.
\end{theorem}

\begin{proof}
Coordinatewise application of Theorem \ref{thm:chain-cube} commutes with restriction maps because all restrictions are coordinate projections.  Hence the maps are sheaf morphisms.  The stalkwise adjunctions are exactly Theorem \ref{thm:chain-cube} at the multiplicity $m_x$.
\end{proof}

\subsection{Derivative-enhanced stalks}

\begin{definition}
Let $f\in\cO(X)$ be a nonzero holomorphic function with zero divisor
\[
\operatorname{div}(f)=D=\sum_{x\in Z}m_x[x].
\]
For each $x\in Z$, set
\[
Q_x:=\cO_{X,x}/\mathfrak m_x^{m_x},
\qquad
K_x:=\cO_{X,x}/\mathfrak m_x^{m_x+1},
\]
and write $[f]_x$ for the class of $f$ in $K_x$.
Define the sheaves
\[
\cQ_D(U):=\prod_{x\in U\cap Z}Q_x,
\qquad
\cK_D(U):=\prod_{x\in U\cap Z}K_x,
\]
and the distinguished section
\[
[f]_D(U):=([f]_x)_{x\in U\cap Z}\in \cK_D(U).
\]
\end{definition}

\begin{proposition}\label{prop:derivative-enhanced-local}
For each $x\in Z$, the class $[f]_x$ is nonzero and lies in the minimal nonzero ideal
\[
\mathfrak m_x^{m_x}/\mathfrak m_x^{m_x+1}\subset K_x.
\]
Indeed, $([f]_x)$ is exactly that minimal nonzero ideal.
\end{proposition}

\begin{proof}
By definition of multiplicity, $\ord_x(f)=m_x$.  Hence in the discrete valuation ring $\cO_{X,x}$ one may write
\[
f=u_x\pi_x^{m_x}
\]
with $u_x$ a unit and $\pi_x$ a uniformizer.  Modulo $\mathfrak m_x^{m_x+1}$ this gives
\[
[f]_x=[u_x]\,[\pi_x]^{m_x},
\]
which is nonzero because $u_x$ remains invertible and $\pi_x^{m_x}\notin\mathfrak m_x^{m_x+1}$.  It lies in $\mathfrak m_x^{m_x}/\mathfrak m_x^{m_x+1}$ by construction.  Since that quotient is one-dimensional over the residue field, every nonzero element generates it, so the ideal generated by $[f]_x$ is exactly the minimal nonzero ideal.
\end{proof}

\begin{corollary}\label{cor:coordinate-derivative}
Let $x\in Z$ and choose a centered local coordinate $t$ at $x$.  Then in $K_x\cong \CC[t]/(t^{m_x+1})$ one has
\[
[f]_x=\lambda_x t^{m_x}
\qquad\text{with}\qquad
\lambda_x\in\CC^\times.
\]
If $X=\CC$ with global coordinate $s$, then
\[
\lambda_x=\frac{f^{(m_x)}(x)}{m_x!}.
\]
\end{corollary}

\begin{proof}
In the chosen coordinate, $f(t)=\lambda_x t^{m_x}+O(t^{m_x+1})$ with $\lambda_x\neq0$ because $\ord_x(f)=m_x$.  Passing to the quotient by $(t^{m_x+1})$ gives the first formula.  On $\CC$, this is exactly Taylor's theorem in the coordinate $s$.
\end{proof}

\begin{remark}
The coordinate-free datum is the minimal nonzero ideal generated by $[f]_x$.  The derivative coefficient $\lambda_x$ appears only after one chooses a local coordinate.  This cleanly separates the intrinsic jet information from its concrete coordinate expression.
\end{remark}

\section{Compactification and tail sheaves at infinity}

We now specialize to $X=\CC$ and compactify by
\[
j:\CC\into\PP^1(\CC),\qquad i_\infty:\{\infty\}\into\PP^1(\CC).
\]
Let $Z\subset\CC$ be a closed discrete subset.

\subsection{General tail objects}

\begin{definition}
Let $(M_x)_{x\in Z}$ be a family of pointed sets with basepoints $*_x$.  Define the sheaf on $\CC$
\[
\mathcal F_M(U):=\prod_{x\in U\cap Z}M_x.
\]
The \emph{tail pointed set} is
\[
\Tail(M):=\left(\prod_{x\in Z}M_x\right)\Big/\sim,
\]
where
\[
(a_x)\sim (b_x)
\iff a_x=b_x\text{ for all but finitely many }x\in Z.
\]
If each $M_x$ carries extra structure preserved by coordinatewise operations, then $\Tail(M)$ inherits the same structure.
\end{definition}

\begin{definition}
For an open set $V\subset\PP^1(\CC)$, define
\[
(j_*\mathcal F_M)(V):=\mathcal F_M(V\cap\CC)=\prod_{x\in V\cap Z}M_x.
\]
If $s=(a_x)\in (j_*\mathcal F_M)(V)$, define its support by
\[
\supp(s):=\{x\in V\cap Z:a_x\neq *_x\}.
\]
Define $j_!\mathcal F_M(V)$ to be the subset of sections whose support is closed in $V$:
\[
(j_!\mathcal F_M)(V):=\{s\in (j_*\mathcal F_M)(V):\supp(s)\text{ is closed in }V\}.
\]
Because $Z$ is discrete in $\CC$, this means:
\begin{enumerate}[label=(\alph*)]
\item if $\infty\notin V$, then $j_!\mathcal F_M(V)=j_*\mathcal F_M(V)$;
\item if $\infty\in V$, then $j_!\mathcal F_M(V)$ consists exactly of finitely supported families on $V\cap Z$.
\end{enumerate}
\end{definition}

\begin{lemma}\label{lem:jshriek-sheaf}
The assignment $j_!\mathcal F_M$ is a sheaf on $\PP^1(\CC)$.
\end{lemma}

\begin{proof}
Let $V=\bigcup_i V_i$ be an open cover and let $s_i\in j_!\mathcal F_M(V_i)$ be compatible.  By Lemma \ref{lem:product-sheaf}, they glue uniquely to a section $s\in j_*\mathcal F_M(V)$.  For each $i$, the support of the restriction of $s$ to $V_i$ is exactly $\supp(s_i)$, which is closed in $V_i$.  Therefore $V_i\setminus\supp(s)$ is open in $V_i$, hence open in $V$.  Since
\[
V\setminus\supp(s)=\bigcup_i (V_i\setminus\supp(s)),
\]
this complement is open in $V$, so $\supp(s)$ is closed in $V$.  Hence $s\in j_!\mathcal F_M(V)$.

The locality statement is immediate: if $s\in j_*\mathcal F_M(V)$ restricts to a section of $j_!\mathcal F_M(V_i)$ on each $V_i$, then the same argument shows that $\supp(s)$ is closed in $V$.
\end{proof}

\begin{theorem}\label{thm:stalks-at-infty}
Let $(M_x)_{x\in Z}$ be a family of pointed sets.  Then the canonical morphism
\[
j_!\mathcal F_M\longrightarrow j_*\mathcal F_M
\]
is an isomorphism on $\PP^1(\CC)\setminus\{\infty\}$, and its stalks satisfy:
\begin{enumerate}[label=(\roman*)]
\item for $x\in Z$, both stalks are $M_x$;
\item for $y\in\CC\setminus Z$, both stalks are one-point stalks;
\item at $\infty$,
\[
(j_!\mathcal F_M)_\infty\cong *,
\qquad
(j_*\mathcal F_M)_\infty\cong \Tail(M).
\]
\end{enumerate}
Consequently, if $\Tail(M)$ is not a singleton, then
\[
j_!\mathcal F_M\not\cong j_*\mathcal F_M
\]
in sheaves of pointed sets.
\end{theorem}

\begin{proof}
For $x\in Z$ or $y\in\CC\setminus Z$, choose a small open neighborhood avoiding all other points of $Z$.  Then both sheaves are either the constant stalk $M_x$ or the one-point stalk.

Now consider $\infty$.  Every germ in $(j_!\mathcal F_M)_\infty$ is represented by a finitely supported section on some neighborhood $V\ni\infty$.  Shrinking $V$ to a smaller neighborhood excluding that finite support turns the section into the basepoint section.  Hence the stalk of $j_!\mathcal F_M$ at $\infty$ is the one-point stalk.

A germ in $(j_*\mathcal F_M)_\infty$ is represented by a section on some neighborhood $V\ni\infty$, i.e. by a family on the cofinite subset $V\cap Z\subset Z$.  Extending by basepoints on the finite complement yields a full family $(a_x)_{x\in Z}$.  Two such families define the same germ exactly when they agree on some smaller neighborhood of $\infty$, i.e. for all but finitely many $x\in Z$.  This is precisely the equivalence relation defining $\Tail(M)$.

If $\Tail(M)$ is nontrivial, the stalks at $\infty$ differ, so the sheaves cannot be isomorphic.
\end{proof}

\begin{definition}
If $(V_x)_{x\in Z}$ is a family of vector spaces, define the tail vector space
\[
\Tail_{\mathrm{vec}}(V):=\left(\prod_{x\in Z}V_x\right)\Big/\left(\bigoplus_{x\in Z}V_x\right).
\]
\end{definition}

\begin{proposition}\label{prop:vector-tail-exact}
Let $(V_x)_{x\in Z}$ be a family of complex vector spaces, and let $\mathcal F_V$ be the corresponding sheaf on $\CC$.  Then there is a short exact sequence of sheaves of vector spaces on $\PP^1(\CC)$,
\[
0\longrightarrow j_!\mathcal F_V
\longrightarrow j_*\mathcal F_V
\longrightarrow i_{\infty*}\Tail_{\mathrm{vec}}(V)
\longrightarrow 0.
\]
Moreover, $j_*\mathcal F_V$ and $i_{\infty*}\Tail_{\mathrm{vec}}(V)$ are flabby, and therefore
\[
H^q\bigl(\PP^1(\CC),j_!\mathcal F_V\bigr)=0\qquad(q\ge1).
\]
Finally,
\[
\Tail_{\mathrm{vec}}(V)
\cong
\Gamma\bigl(\PP^1(\CC),j_*\mathcal F_V\bigr)\Big/\Gamma\bigl(\PP^1(\CC),j_!\mathcal F_V\bigr).
\]
\end{proposition}

\begin{proof}
Away from $\infty$, the two sheaves agree, so the quotient sheaf is supported at $\infty$.  On a neighborhood $V\ni\infty$, sections of $j_*\mathcal F_V$ are arbitrary families on $V\cap Z$, while sections of $j_!\mathcal F_V$ are finitely supported families.  Since $V\cap Z$ is cofinite in $Z$, extension by zero on the finite complement identifies the quotient with $\Tail_{\mathrm{vec}}(V)$, independently of $V$.  This gives the short exact sequence.

Restriction maps of $j_*\mathcal F_V$ are coordinate projections between products, hence surjective; so $j_*\mathcal F_V$ is flabby.  The skyscraper sheaf $i_{\infty*}\Tail_{\mathrm{vec}}(V)$ is also flabby.  The long exact cohomology sequence then gives $H^q(\PP^1,j_!\mathcal F_V)=0$ for $q\ge2$ and an exact segment
\[
0\to \Gamma(j_!\mathcal F_V)\to \Gamma(j_*\mathcal F_V)\to \Tail_{\mathrm{vec}}(V)
\to H^1(\PP^1,j_!\mathcal F_V)\to 0.
\]
But the map on global sections is by definition the quotient map from the product to the direct-sum quotient, so it is surjective.  Hence $H^1$ also vanishes, and the quotient description follows.
\end{proof}

\subsection{Boolean tails}

\begin{definition}
For a divisor $D=\sum_{x\in Z}m_x[x]$ on $\CC$, define the Boolean tail semimodules
\[
T_D^{\cV}:=\Tail\bigl((A^{m_x})_{x\in Z}\bigr),
\qquad
T_D^{\cA}:=\Tail\bigl((A)_{x\in Z}\bigr).
\]
\end{definition}

\begin{corollary}\label{cor:boolean-tail-sheaf}
For the Boolean shadow sheaves of a divisor on $\CC$, the canonical morphisms
\[
j_!\cV_D\to j_*\cV_D,
\qquad
j_!\cA_D\to j_*\cA_D
\]
are isomorphisms away from $\infty$, while the stalks at $\infty$ are
\[
(j_!\cV_D)_\infty\cong 0,
\qquad
(j_*\cV_D)_\infty\cong T_D^{\cV},
\]
\[
(j_!\cA_D)_\infty\cong 0,
\qquad
(j_*\cA_D)_\infty\cong T_D^{\cA}.
\]
If $D$ has infinite support, then both tail semimodules are nonzero.
\end{corollary}

\begin{proof}
Apply Theorem \ref{thm:stalks-at-infty} to the pointed sets underlying the semimodules.  If the support is infinite, the class of the everywhere-nonzero family is not eventually equal to the basepoint family, so the tail object is nontrivial.
\end{proof}

\section{Zeta applications}

\subsection{The zeta divisor and its split-zero shadow}

\begin{definition}
Let
\[
D_\zeta:=\sum_{\rho:\,\zeta(\rho)=0} m_\rho[\rho]
\]
be the effective divisor of zeta zeros on $\CC$.  Let
\[
\cV_\zeta:=\cV_{D_\zeta},\qquad \cA_\zeta:=\cA_{D_\zeta},
\qquad \cI_\zeta:=\cI_{D_\zeta},\qquad \cQ_\zeta:=\cQ_{D_\zeta},\qquad \cK_\zeta:=\cK_{D_\zeta}.
\]
\end{definition}

\begin{corollary}\label{cor:zeta-shadow-sheaves}
At every zero $\rho$ of zeta,
\[
(\cV_\zeta)_\rho\cong A^{m_\rho},
\qquad
(\cA_\zeta)_\rho\cong A,
\qquad
\Fix_S(F_1(\rho))=A.
\]
The reduced zeta shadow is a canonical split retract of the multiplicity shadow.
\end{corollary}

\begin{proof}
Combine Proposition \ref{prop:shadow-stalks}, Proposition \ref{prop:retract-shadow}, and Corollary \ref{cor:zeta-pointwise}.
\end{proof}

\subsection{Completed zeta and packet descent}

\begin{definition}
Define the completed zeta function
\[
\xi(s):=\frac12 s(s-1)\pi^{-s/2}\Gamma\!\left(\frac s2\right)\zeta(s).
\]
Let $D_\xi$ be its zero divisor.  Thus $D_\xi$ is the divisor of nontrivial zeros of zeta.
\end{definition}

\begin{definition}
Let
\[
c(s):=\overline{s},
\qquad
w(s):=1-s,
\qquad
W:=\{\mathrm{id},c,w,cw\}\cong C_2\times C_2.
\]
\end{definition}

\begin{lemma}\label{lem:xi-symmetry-multiplicity}
If $\rho$ is a zero of $\xi$, then $\overline\rho$, $1-\rho$, and $1-\overline\rho$ are also zeros of $\xi$, and all four points occur with the same multiplicity.
\end{lemma}

\begin{proof}
The functional equation $\xi(1-s)=\xi(s)$ gives the reflection symmetry.  The Dirichlet series for $\zeta$ has real coefficients on $\Re(s)>1$, so $\overline{\zeta(\overline s)}=\zeta(s)$ there, and analytic continuation gives the same relation meromorphically on all of $\CC$.  The elementary prefactor in $\xi$ is also conjugation-compatible, so $\overline{\xi(\overline s)}=\xi(s)$.  Therefore zeros are preserved by conjugation.  Multiplicities are preserved because both symmetries are biholomorphic or antiholomorphic automorphisms with nonzero local derivative, so they preserve vanishing order.
\end{proof}

\begin{definition}
Let $Y_\xi:=|D_\xi|/W$ be the orbit space of nontrivial zero packets, equipped with the quotient topology.  Since $|D_\xi|$ is discrete, so is $Y_\xi$.  For an orbit $O\in Y_\xi$, write $m_O$ for the common multiplicity of its points.

Define the \emph{Boolean packet sheaf}
\[
\mathcal P_\xi^{\cV}(U):=\prod_{O\in U} A^{m_O}
\]
and the \emph{chain packet sheaf}
\[
\mathcal P_\xi^{\cC}(U):=\prod_{O\in U} C_{m_O}
\]
on open sets $U\subset Y_\xi$.
\end{definition}

\begin{theorem}\label{thm:packet-descent}
Let $q:|D_\xi|\to Y_\xi$ be the quotient map.  The group $W$ acts naturally on the sheaves $\cV_{D_\xi}$ and $\cC_{D_\xi}$.  For each orbit $O\in Y_\xi$ of cardinality $r_O$, the stalk of $q_*\cV_{D_\xi}$ at $O$ is canonically
\[
(q_*\cV_{D_\xi})_O\cong (A^{m_O})^{r_O},
\]
and its $W$-fixed subsemimodule is the diagonal copy of $A^{m_O}$.  Hence
\[
(q_*\cV_{D_\xi})^W\cong \mathcal P_\xi^{\cV}.
\]
Likewise,
\[
(q_*\cC_{D_\xi})_O\cong (C_{m_O})^{r_O},
\qquad
(q_*\cC_{D_\xi})^W\cong \mathcal P_\xi^{\cC}.
\]
Thus both the Boolean multiplicity shadow and the classical jet-ideal chain descend to honest packet sheaves on the orbit space.
\end{theorem}

\begin{proof}
Because $Y_\xi$ is discrete, the stalk of $q_*\cV_{D_\xi}$ at an orbit $O$ is simply the product of the stalks over the points in that orbit.  Lemma \ref{lem:xi-symmetry-multiplicity} gives the common multiplicity $m_O$, so
\[
(q_*\cV_{D_\xi})_O\cong \prod_{\rho\in O} A^{m_O}\cong (A^{m_O})^{r_O}.
\]
The group $W$ acts by permuting the $r_O$ factors transitively.  An element of $(A^{m_O})^{r_O}$ is invariant under this permutation action if and only if all its coordinates are equal.  Thus the invariant subsemimodule is the diagonal copy of $A^{m_O}$.  Since the orbit space is discrete, these diagonal stalks assemble into the sheaf $\mathcal P_\xi^{\cV}$.

The chain case is identical, replacing $A^{m_O}$ by $C_{m_O}$.
\end{proof}

\begin{remark}
The packet sheaf is the precise sheaf-theoretic version of the idea that nontrivial zeta zeros should be studied in symmetry packets rather than as isolated points.  No hypothesis such as RH is required.  If RH holds, the orbit sizes are $2$; off the critical line, the orbit sizes are $4$.
\end{remark}

\subsection{Trivial zeros and exact stalk weights}

\begin{proposition}\label{prop:trivial-simple}
For every integer $n\ge1$, the trivial zero $-2n$ of zeta is simple.
\end{proposition}

\begin{proof}
The functional equation reads
\[
\zeta(s)=2^s\pi^{s-1}\sin\!\left(\frac{\pi s}{2}\right)\Gamma(1-s)\zeta(1-s).
\]
At $s=-2n$, the factors $2^s$, $\pi^{s-1}$, and $\Gamma(1-s)=\Gamma(1+2n)$ are holomorphic and nonzero, while $1-s=1+2n>1$ implies $\zeta(1-s)=\zeta(2n+1)\neq0$.  The only vanishing factor is $\sin(\pi s/2)$, and it has a simple zero because
\[
\frac{d}{ds}\sin\!\left(\frac{\pi s}{2}\right)\Big|_{s=-2n}=
\frac\pi2\cos(-n\pi)=\frac\pi2(-1)^n\neq0.
\]
Therefore $-2n$ is a simple zero.
\end{proof}

\begin{corollary}\label{cor:trivial-stalk}
On the trivial-zero sector, the Boolean multiplicity stalk is just $A$:
\[
(\cV_\zeta)_{-2n}\cong A\qquad(n\ge1).
\]
\end{corollary}

\begin{proof}
This is Proposition \ref{prop:shadow-stalks} together with Proposition \ref{prop:trivial-simple}.
\end{proof}

\begin{proposition}\label{prop:trivial-derivative}
For every integer $n\ge1$,
\[
\zeta'(-2n)=(-1)^n\frac{(2n)!}{2(2\pi)^{2n}}\,\zeta(2n+1).
\]
Hence
\[
|\zeta'(-2n)|=\frac{(2n)!}{2(2\pi)^{2n}}\,\zeta(2n+1).
\]
\end{proposition}

\begin{proof}
Set
\[
g_n(s):=2^s\pi^{s-1}\Gamma(1-s)\zeta(1-s).
\]
Then $\zeta(s)=g_n(s)\sin(\pi s/2)$ and $g_n$ is holomorphic at $s=-2n$ with
\[
g_n(-2n)=2^{-2n}\pi^{-2n-1}(2n)!\,\zeta(2n+1).
\]
Differentiating and using that the sine factor vanishes at $-2n$ gives
\[
\zeta'(-2n)=g_n(-2n)\cdot \frac\pi2\cos(-n\pi)
=2^{-2n}\pi^{-2n-1}(2n)!\,\zeta(2n+1)\cdot\frac\pi2(-1)^n.
\]
This simplifies to the displayed formula.  Because $\zeta(2n+1)>0$, the absolute-value formula follows.
\end{proof}

\begin{definition}
For each trivial zero $-2n$, define its \emph{canonical Boolean stalk weight}
\[
\omega_{-2n}:A\to\RR_{\ge0}
\]
by
\[
\omega_{-2n}(\taup):=0,
\qquad
\omega_{-2n}(0):=|\zeta'(-2n)|.
\]
\end{definition}

\begin{corollary}\label{cor:trivial-weight}
The unique nonempty face of the trivial Boolean stalk carries the exact analytic weight
\[
\omega_{-2n}(0)=\frac{(2n)!}{2(2\pi)^{2n}}\,\zeta(2n+1).
\]
\end{corollary}

\begin{proof}
This is immediate from the definition and Proposition \ref{prop:trivial-derivative}.
\end{proof}

\subsection{Compactified zeta tails}

\begin{definition}
Define the zeta tail objects
\[
T_\zeta^{\cV}:=\Tail\bigl((A^{m_\rho})_{\rho:\,\zeta(\rho)=0}\bigr),
\qquad
T_\zeta^{\cA}:=\Tail\bigl((A)_{\rho:\,\zeta(\rho)=0}\bigr).
\]
Define also the jet-vector-space tails
\[
T_\zeta^{\cQ}:=\Tail_{\mathrm{vec}}\bigl((Q_\rho)_{\rho:\,\zeta(\rho)=0}\bigr),
\qquad
T_\zeta^{\cK}:=\Tail_{\mathrm{vec}}\bigl((K_\rho)_{\rho:\,\zeta(\rho)=0}\bigr).
\]
\end{definition}

\begin{theorem}\label{thm:zeta-tail}
The compactified Boolean and jet shadows of zeta satisfy:
\begin{enumerate}[label=(\roman*)]
\item the stalks at $\infty$ are
\[
(j_*\cV_\zeta)_\infty\cong T_\zeta^{\cV},
\qquad
(j_*\cA_\zeta)_\infty\cong T_\zeta^{\cA},
\]
\[
(j_*\cQ_\zeta)_\infty\cong T_\zeta^{\cQ},
\qquad
(j_*\cK_\zeta)_\infty\cong T_\zeta^{\cK};
\]
\item all four tail objects are nonzero;
\item for the jet-vector-space shadows, there are short exact sequences
\[
0\to j_!\cQ_\zeta\to j_*\cQ_\zeta\to i_{\infty*}T_\zeta^{\cQ}\to0,
\]
\[
0\to j_!\cK_\zeta\to j_*\cK_\zeta\to i_{\infty*}T_\zeta^{\cK}\to0,
\]
and therefore
\[
H^q\bigl(\PP^1(\CC),j_!\cQ_\zeta\bigr)=H^q\bigl(\PP^1(\CC),j_!\cK_\zeta\bigr)=0
\qquad(q\ge1).
\]
\end{enumerate}
\end{theorem}

\begin{proof}
Part (i) is Theorem \ref{thm:stalks-at-infty} for Boolean shadows and Proposition \ref{prop:vector-tail-exact} for jet shadows.  Since zeta has infinitely many trivial zeros, the index set is infinite; the class of the everywhere-nonbasepoint family on the trivial-zero sector shows that the Boolean tails are nonzero, and any everywhere-nonzero vector in the product gives nonzero jet tails.  This proves (ii).  Part (iii) is Proposition \ref{prop:vector-tail-exact}.
\end{proof}

\section{Conclusion}

The pointwise split-zero fixed locus of a zeta zero is the Boolean pair $\{0,\taup\}$.  The mathematical content begins only after one organizes all zeros at once.

This paper has produced four such organizations.

\begin{enumerate}[label=(\roman*)]
\item The Boolean multiplicity shadow $\cV_D$, whose stalk at multiplicity $m$ is the Boolean cube $B_m$.
\item The reduced Boolean shadow $\cA_D$, which is a canonical split retract of $\cV_D$.
\item The classical local Artin ring $Q_{x,m}$ and its ideal lattice $C_m$, which provide the genuine coordinate-free nonreduced thickening.
\item The compactified tail sheaf at $\infty$, which is the first genuinely global invariant visible in the present discrete setting.
\end{enumerate}

The decisive comparison theorem is not an isomorphism but a biadjoint relation:
\[
L_m\dashv i_m\dashv R_m
\]
between the jet-ideal chain and the Boolean cube.  This gives a precise, category-correct relation between the classical nonreduced divisor and its split-zero Boolean shadow.

On the analytic side, the local jet class of a holomorphic function with divisor $D$ generates the minimal nonzero ideal in the $(m+1)$-jet algebra.  For zeta this recovers the exact derivatives at the zeros; on the trivial sector it yields the explicit weight
\[
|\zeta'(-2n)|=\frac{(2n)!}{2(2\pi)^{2n}}\,\zeta(2n+1).
\]
On the symmetry side, the nontrivial zeros of the completed zeta function descend to genuine packet sheaves on the Klein-four orbit space.

The outcome is therefore a complete theorem package, not a slogan: the zeta zero divisor has a Boolean multiplicity shadow, a classical jet-ideal chain, a coordinate-free derivative-enhanced minimal ideal, a compactified tail at infinity, and a symmetry-packet descent.  Each object is explicit, and the relations among them are proved.

\begin{thebibliography}{99}

\bibitem{DrosteKuichVogler}
M. Droste, W. Kuich, and H. Vogler (eds.),
\emph{Handbook of Weighted Automata},
Springer, Berlin, 2009.

\bibitem{GaubertKatz}
S. Gaubert and R. Katz,
\emph{Rational semimodules over the max-plus semiring and geometric approach of discrete event systems},
arXiv:math/0208014, 2002.

\bibitem{GrothendieckEGAII}
A. Grothendieck,
\emph{\'El\'ements de g\'eom\'etrie alg\'ebrique. II. \`Etude globale \`el\'ementaire de quelques classes de morphismes},
Publ. Math. IH\'ES \textbf{8} (1961), 5--222.

\bibitem{Gunawardena}
J. Gunawardena (ed.),
\emph{Idempotency},
Cambridge University Press, Cambridge, 1998.

\bibitem{JarraLorscheidVital}
M. Jarra, O. Lorscheid, and E. Vital,
\emph{Quiver matroids, matroid morphisms, quiver Grassmannians, their Euler characteristics and $\mathbb F_1$-points},
arXiv:2404.09255, version 2, 2026.

\bibitem{Lee}
J. M. Lee,
\emph{Introduction to Complex Manifolds},
American Mathematical Society, Providence, RI, 2024.

\bibitem{Lorscheid}
O. Lorscheid,
\emph{The geometry of blueprints. Part I: Algebraic background and scheme theory},
Adv. Math. \textbf{229} (2012), no. 3, 1804--1846.

\bibitem{OrecchiaChiantini}
F. Orecchia and L. Chiantini (eds.),
\emph{Zero-Dimensional Schemes},
Walter de Gruyter, Berlin, 1994.

\bibitem{Vakil}
R. Vakil,
\emph{The Rising Sea: Foundations of Algebraic Geometry},
July 27, 2024 draft.

\end{thebibliography}

\end{document}
